% © 2026 Ing. David Jaroš — CC BY-NC-ND 4.0
%
% UBT-AI-PROVENANCE-BEGIN
% schema: ubt-ai-provenance/v1
% tier: C_working
% ai_assistance: disclosed
% human_review: risk-based
% editorial_responsibility: Ing. David Jaroš
% policy: ../../AI_PROVENANCE.md
% notice: Working material; exhaustive human review is not claimed.
% UBT-AI-PROVENANCE-END

\ifdefined\INCLUDEMODE
\else
  \documentclass[12pt,a4paper]{article}
  \usepackage[a4paper, margin=2.5cm]{geometry}
  \usepackage{amsmath, amssymb, amsthm}
  \usepackage{mathtools}
  \usepackage{hyperref}
  \usepackage{xcolor}
  \usepackage{mdframed}
  \usepackage{booktabs}

  \newtheorem{theorem}{Theorem}[section]
  \newtheorem{lemma}[theorem]{Lemma}
  \newtheorem{proposition}[theorem]{Proposition}
  \newtheorem{corollary}[theorem]{Corollary}
  \theoremstyle{definition}
  \newtheorem{definition}[theorem]{Definition}
  \theoremstyle{remark}
  \newtheorem{remark}[theorem]{Remark}

  \newmdenv[backgroundcolor=yellow!10, linecolor=orange!80!black, linewidth=1.5pt]{gapbox}
  \newmdenv[backgroundcolor=green!8, linecolor=green!60!black, linewidth=1.5pt]{resultbox}
  \newmdenv[backgroundcolor=blue!5, linecolor=blue!50, linewidth=1.5pt]{statusbox}

  \title{\textbf{Canonical UBT Colour Sector: \\
         $\mathrm{SU}(3)$ Stabilizer and Exterior-Fock Derivation}}
  \author{Ing.\ David Jaro\v{s}}
  \date{2026}

  \begin{document}
  \maketitle

\begin{abstract}
Starting exclusively from the biquaternionic algebra $\mathbb{C}\otimes\mathbb{H}$
already present in UBT, we derive a canonical Hermitian form $h$ and a canonical
complex volume form $\Omega$ on the colour subspace
$V = \mathbb{C}\text{-span}\{I,J,K\}$, both obtained from the biquaternionic
adjoint and quaternion multiplication.  We prove that the joint stabiliser of
$h$ and $\Omega$ in $\mathrm{GL}(3,\mathbb{C})$ is exactly $\mathrm{SU}(3)$.
We construct the exterior algebra $\Lambda^\bullet V \cong \mathbf{1}\oplus\mathbf{3}
\oplus\bar{\mathbf{3}}\oplus\mathbf{1}$, show that the $8$-dimensional
second-quantised representation of $\mathfrak{su}(3)$ acts on this space, and
verify the complete commutator algebra exactly using SymPy.
The proof is self-contained within $\mathbb{C}\otimes\mathbb{H}$; it does not
introduce octonions, fibres, or any new algebraic structure.
The dynamical identification of this $\mathrm{SU}(3)$ with the local QCD
gauge group is explicitly left as \textbf{GAP-SU3-DYN: OPEN}.
\end{abstract}
\fi

%──────────────────────────────────────────────────────────────────────────────
\section{Status and claim discipline}
\label{sec:status}
%──────────────────────────────────────────────────────────────────────────────

\begin{statusbox}
\begin{tabular}{@{}p{0.56\linewidth}p{0.15\linewidth}p{0.23\linewidth}@{}}
\toprule
Item & Status & Notes \\
\midrule
Canonical colour space $V=\mathbb{C}\text{-span}\{I,J,K\}\cong\mathbb{C}^3$
  & \textbf{Proved} [L1] & \S\ref{sec:colour_space}\\
Hermitian form $h$ from biquaternionic $\dagger$
  & \textbf{Proved} [L1] & \S\ref{sec:hermitian}\\
Volume form $\Omega$ from quaternion multiplication
  & \textbf{Proved} [L1] & \S\ref{sec:volume}\\
$\mathrm{Stab}(h,\Omega)=\mathrm{SU}(3)$ — stabiliser theorem
  & \textbf{Proved} [L1] & \S\ref{sec:stabiliser}\\
Correct global structure $U(3)\cong(SU(3)\times U(1))/\mathbb{Z}_3$
  & \textbf{Stated} & \S\ref{sec:global_group}\\
$\Lambda^\bullet V\cong\mathbf{1}\oplus\mathbf{3}\oplus\bar{\mathbf{3}}\oplus\mathbf{1}$
  & \textbf{Proved} [L1] & \S\ref{sec:exterior}\\
$8$ second-quantised $\mathfrak{su}(3)$ generators on $\Lambda^\bullet V$
  & \textbf{Proved} [L1] & \S\ref{sec:fock}\\
Machine verification (SymPy, exact)
  & \textbf{PASS} & \texttt{verification/su3\_stabilizer\_fock\_check.py}\\
Lean formalisation
  & \textbf{LEAN-PENDING} & full formalisation in progress\\
\midrule
Local dynamical identification with QCD gauge field
  & \textbf{OPEN} & GAP-SU3-DYN\\
Dynamical confinement / mass gap
  & \textbf{OPEN} & Clay Millennium problem\\
\bottomrule
\end{tabular}
\end{statusbox}

%──────────────────────────────────────────────────────────────────────────────
\section{Historical context}
\label{sec:history}
%──────────────────────────────────────────────────────────────────────────────

Earlier UBT approaches to $\mathrm{SU}(3)_c$ are preserved in
\texttt{canonical/su3\_derivation/archive/}.  Specifically:
\begin{itemize}
  \item The \emph{involution approach} (Theorems G.A--G.D in
        \texttt{su3\_from\_involutions.tex}) proved the carrier selection
        $V_c\cong\mathbb{C}^3$ but \emph{postulated} the unitary structure on $V_c$.
  \item The \emph{superposition/triqubit approach} correctly identified the carrier
        but used an incorrect global-group statement
        $U(3)\cong U(1)\times\mathrm{SU}(3)$; see \S\ref{sec:global_group}.
  \item Error-correction and I-Ching-based motivations (triqubit QEC) appear in
        historical files.  They are preserved for provenance but play no role
        in the canonical proof.
\end{itemize}

The present derivation differs in that it obtains the Hermitian form and volume
form \emph{canonically} from biquaternionic operations already present in UBT,
without any postulate about unitarity.

%──────────────────────────────────────────────────────────────────────────────
\section{The canonical colour space}
\label{sec:colour_space}
%──────────────────────────────────────────────────────────────────────────────

The UBT field takes values in $\mathbb{C}\otimes_{\mathbb{R}}\mathbb{H}$.
As a $\mathbb{C}$-vector space this algebra has complex basis
$\{1,I,J,K\}$ and hence complex dimension $4$; equivalently, as a
$\mathbb{R}$-vector space it has real basis
$\{1,I,J,K,i,iI,iJ,iK\}$ and real dimension $8$.  Here $1$ is the quaternion
identity and $I,J,K$ are the three imaginary quaternion units satisfying
\begin{equation}\label{eq:quat}
  I^2=J^2=K^2=-1,\quad IJ=K,\quad JK=I,\quad KI=J.
\end{equation}
We define the \emph{canonical colour subspace}
\begin{equation}\label{eq:V_def}
  V := \mathbb{C}\text{-span}\{I,J,K\} \subset \mathbb{C}\otimes\mathbb{H}.
\end{equation}
Since $\{I,J,K\}$ are a subset of the complex basis $\{1,I,J,K\}$,
they are $\mathbb{C}$-linearly independent and $V$ has complex dimension $3$.

\begin{proposition}\label{prop:V_iso}
$V \cong \mathbb{C}^3$ as a $\mathbb{C}$-vector space via
$\alpha I+\beta J+\gamma K \leftrightarrow (\alpha,\beta,\gamma)$.
\end{proposition}

This isomorphism is the unique one mapping the ordered basis $(I,J,K)$ to the
standard basis of $\mathbb{C}^3$.  An arbitrary element of $V$ is written
$v = v_1 I + v_2 J + v_3 K$ with $v_1,v_2,v_3\in\mathbb{C}$.

%──────────────────────────────────────────────────────────────────────────────
\section{Hermitian form from the biquaternionic adjoint}
\label{sec:hermitian}
%──────────────────────────────────────────────────────────────────────────────

\subsection{The biquaternionic dagger}

The UBT biquaternionic adjoint $(\cdot)^\dagger$ is the composite
\begin{equation}\label{eq:dagger}
  (z\otimes q)^\dagger := \bar{z}\otimes\bar{q},
\end{equation}
where $\bar{z}$ is complex conjugation and $\bar{q}$ is quaternionic
conjugation ($\overline{q_0+q_i I+q_j J+q_k K}=q_0-q_i I-q_j J-q_k K$).
On the colour subspace, for $v = v_1 I+v_2 J+v_3 K\in V$,
\begin{equation}
  v^\dagger = \bar{v}_1(-I)+\bar{v}_2(-J)+\bar{v}_3(-K)
            = -\bar{v}_1 I -\bar{v}_2 J -\bar{v}_3 K.
\end{equation}

\subsection{The scalar projection}

The \emph{scalar part} of a biquaternion $x = x_0 + x_I I + x_J J + x_K K$
(with $x_0\in\mathbb{C}$) is $\mathrm{Sc}(x) := x_0$.

\subsection{Canonical Hermitian form}

\begin{definition}\label{def:h}
For $v,w\in V$, define
\begin{equation}\label{eq:h_def}
  h(v,w) := \mathrm{Sc}(v^\dagger w).
\end{equation}
\end{definition}

\begin{theorem}\label{thm:hermitian}
The form $h$ defined in \eqref{eq:h_def} satisfies, for
$v=v_1 I+v_2 J+v_3 K$ and $w=w_1 I+w_2 J+w_3 K$:
\begin{equation}\label{eq:h_explicit}
  h(v,w) = \bar{v}_1 w_1 + \bar{v}_2 w_2 + \bar{v}_3 w_3.
\end{equation}
In particular, $h$ is the standard Hermitian inner product on $\mathbb{C}^3$.
\end{theorem}

\begin{proof}
Using $I^2=J^2=K^2=-1$ and the fact that the products of distinct imaginary
quaternion units have zero scalar part,
\begin{align}
  \mathrm{Sc}(v^\dagger w)
  &= \mathrm{Sc}\!\left[
     (-\bar v_1 I-\bar v_2 J-\bar v_3 K)
     (w_1 I+w_2 J+w_3 K)\right] \\
  &= \bar v_1 w_1+\bar v_2 w_2+\bar v_3 w_3.
\end{align}
Therefore $h(v,w)=\mathrm{Sc}(v^\dagger w)$ is exactly the standard
positive-definite Hermitian product on $\mathbb{C}^3$ in the basis $(I,J,K)$.
\end{proof}

\begin{corollary}
$h$ is:
\begin{enumerate}
  \item \emph{complex sesquilinear}: $h(\lambda v,w)=\bar\lambda h(v,w)$,
        $h(v,\lambda w)=\lambda h(v,w)$;
  \item \emph{Hermitian}: $h(v,w)=\overline{h(w,v)}$;
  \item \emph{positive definite}: $h(v,v)=|v_1|^2+|v_2|^2+|v_3|^2\geq 0$,
        with equality only for $v=0$;
  \item \emph{orthonormal}: $h(I,I)=h(J,J)=h(K,K)=1$,
        $h(I,J)=h(I,K)=h(J,K)=0$.
\end{enumerate}
The group preserving $h$ is therefore $U(3)$.
\end{corollary}

%──────────────────────────────────────────────────────────────────────────────
\section{Volume form from quaternion multiplication}
\label{sec:volume}
%──────────────────────────────────────────────────────────────────────────────

\begin{definition}\label{def:Omega}
For $v,w,u\in V$, define
\begin{equation}\label{eq:Omega_def}
  \Omega(v,w,u) := -\mathrm{Sc}(v\,w\,u).
\end{equation}
\end{definition}

\begin{theorem}\label{thm:Omega}
For $v=v_1 I+v_2 J+v_3 K$, $w=w_1 I+w_2 J+w_3 K$, $u=u_1 I+u_2 J+u_3 K$,
\begin{equation}\label{eq:Omega_det}
  \Omega(v,w,u)
  = \det\begin{pmatrix} v_1 & v_2 & v_3 \\ w_1 & w_2 & w_3 \\ u_1 & u_2 & u_3 \end{pmatrix}.
\end{equation}
In particular $\Omega(I,J,K)=1$.
\end{theorem}

\begin{proof}
Expand the triple product using the quaternion algebra \eqref{eq:quat}.
The scalar part of $XY$ for distinct imaginary units $X,Y\in\{I,J,K\}$ is zero.
The scalar part of $(XY)Z$ for the cyclic product $(I)(J)(K)$ equals the scalar
part of $K\cdot K = -1$, giving $-1$; the sign in $\Omega$ then gives $+1$.
For the anti-cyclic product $(I)(K)(J)$ the scalar part of $(IK)J=(-J)J=1$, so
$-\mathrm{Sc}(IKJ)=-1$, consistent with the determinant being $-1$ for a transposition.

For general coefficients, expand:
\begin{align}
  vwu &= (v_1 I + v_2 J + v_3 K)(w_1 I + w_2 J + w_3 K)(u_1 I + u_2 J + u_3 K).
\end{align}
In the triple product, a term $(\alpha X)(\beta Y)(\gamma Z)$ contributes to the
scalar part only when $XYZ\in\{-1\}$, i.e.\ $XYZ$ is a scalar multiple of the
identity.  Consulting the quaternion multiplication table:
\begin{align*}
  IJK &= K^2 = -1,       & IKJ &= -J^2 = 1, \\
  JIK &= -K^2 = 1,       & JKI &= I^2  = -1, \\
  KIJ &= J^2  = -1,      & KJI &= -I^2 = 1.
\end{align*}
(All other products of three distinct imaginary units include $IJ$, $JK$, $KI$
which land in imaginary, and same-unit pairs $X^2=-1$ are accounted for above.)
Also $XXX = X(XX) = X(-1) = -X$ (imaginary), so pure-repeated triples do not
contribute scalar parts.  For two equal units $XYX$ or similar, they are also
imaginary.  Therefore the scalar part of $vwu$ collects only the six terms with
all three distinct units:
\begin{align}
  \mathrm{Sc}(vwu) &= v_1 w_2 u_3 \cdot(-1) + v_1 w_3 u_2\cdot(+1)
                  + v_2 w_1 u_3\cdot(+1) + v_2 w_3 u_1\cdot(-1) \notag\\
                 &\quad + v_3 w_1 u_2\cdot(-1) + v_3 w_2 u_1\cdot(+1).
\end{align}
Hence
\begin{equation}
  -\mathrm{Sc}(vwu) = v_1 w_2 u_3 - v_1 w_3 u_2 - v_2 w_1 u_3 + v_2 w_3 u_1
                     + v_3 w_1 u_2 - v_3 w_2 u_1,
\end{equation}
which is exactly the Leibniz expansion of
$\det\bigl[(v_i),(w_i),(u_i)\bigr]$.
\end{proof}

\begin{remark}[Orientation]\label{rem:orientation}
The formula $\Omega(v,w,u)=-\mathrm{Sc}(vwu)$ is basis-free once the
biquaternion multiplication and scalar projection are fixed.  In any positively
oriented orthonormal imaginary-quaternion frame $(I,J,K)$ satisfying $IJ=K$ it
has the normalisation $\Omega(I,J,K)=+1$.  Reversing the convention changes only
the overall sign of $\Omega$ and not its stabiliser.
\end{remark}

\begin{corollary}
$\Omega$ is complex trilinear, alternating, and nonzero.  It is the unique
(up to scalar multiple) alternating trilinear form on $V$, and it is normalised
by $\Omega(I,J,K)=1$.
\end{corollary}

%──────────────────────────────────────────────────────────────────────────────
\section{The stabiliser theorem}
\label{sec:stabiliser}
%──────────────────────────────────────────────────────────────────────────────

\begin{theorem}[Canonical colour stabiliser]\label{thm:SU3}
Let $V = \mathbb{C}\text{-span}\{I,J,K\}$ with the canonically induced
Hermitian form $h$ \eqref{eq:h_def} and complex volume form $\Omega$
\eqref{eq:Omega_def}.  Then
\begin{equation}
  \mathrm{Stab}_{\mathrm{GL}(V)}(h,\Omega)
  = \{g\in\mathrm{GL}(3,\mathbb{C}) : g^\dagger g = \mathbf{1},\;\det g = 1\}
  \cong \mathrm{SU}(3).
\end{equation}
\end{theorem}

\begin{proof}
\emph{Step 1 (preservation of $h$ implies $g\in U(3)$).}
If $h(gv,gw)=h(v,w)$ for all $v,w\in V$ then $g^\dagger g = \mathbf{1}_3$,
i.e.\ $g\in U(3)$.

\emph{Step 2 (action of $g$ on $\Omega$).}
For any $g\in\mathrm{GL}(V)$,
\begin{equation}
  \Omega(gv,gw,gu) = \det(g)\,\Omega(v,w,u),
\end{equation}
since $\Omega$ is the determinant form (Theorem~\ref{thm:Omega}).

\emph{Step 3 (preservation of $\Omega$ implies $\det g = 1$).}
If $\Omega(gv,gw,gu)=\Omega(v,w,u)$ for all $v,w,u\in V$ then
$\det(g)=1$, i.e.\ $g\in\mathrm{SL}(3,\mathbb{C})$.

\emph{Step 4 (intersection).}
The common stabiliser is $U(3)\cap\mathrm{SL}(3,\mathbb{C}) = \mathrm{SU}(3)$.

\emph{Step 5 (converse).}
Every $g\in\mathrm{SU}(3)$ satisfies $g^\dagger g=\mathbf{1}$ (preserves $h$)
and $\det g=1$ (preserves $\Omega$).
\end{proof}

\begin{remark}\label{rem:su3_not_automorphism}
$\mathrm{SU}(3)$ is the stabiliser of the \emph{pair} $(h,\Omega)$; it is
not the automorphism group of the biquaternion algebra
$\mathbb{C}\otimes\mathbb{H}$ itself.  Indeed,
$\mathbb{C}\otimes_{\mathbb{R}}\mathbb{H}\cong M_2(\mathbb{C})$, whose
$\mathbb{C}$-algebra automorphisms are inner (isomorphic to
$\mathrm{PGL}_2(\mathbb{C})$ by Skolem--Noether).
\end{remark}

%──────────────────────────────────────────────────────────────────────────────
\section{Correct global group structure}
\label{sec:global_group}
%──────────────────────────────────────────────────────────────────────────────

The group $U(3)$ is the group preserving $h$ alone.
Globally,
\begin{equation}\label{eq:U3_global}
  U(3) \cong \bigl(\mathrm{SU}(3)\times U(1)\bigr)/\mathbb{Z}_3,
\end{equation}
where the $\mathbb{Z}_3$ identification is
$(g,z)\sim(\omega g,\omega^{-1}z)$ with
$\omega=e^{2\pi i/3}$, i.e. the kernel of
$(g,z)\mapsto zg$ is $\{(\omega^k\mathbf 1,\omega^{-k}):k=0,1,2\}$.
This is \emph{not} a direct product.

The Lie-algebra decomposition
\begin{equation}
  \mathfrak{u}(3) = \mathfrak{su}(3)\oplus\mathfrak{u}(1)
\end{equation}
is valid (as a vector-space and even as a Lie-algebra direct sum, since the
centre $\mathfrak{u}(1)$ is abelian), but it does \emph{not} imply a global
direct-product structure.

\begin{remark}[Correction of earlier UBT statement]
The statement $U(3)\cong U(1)\times\mathrm{SU}(3)$ appearing in earlier
versions of \texttt{step1\_superposition\_approach.tex} (Theorem~3.1 there) is
incorrect at the global group level.  The Lie-algebra split is valid; the global
correction is \eqref{eq:U3_global}.  The physical consequence is that the
correct structure for the colour $+$ baryon-number factor is
$(SU(3)\times U(1))/\mathbb{Z}_3$, and care is needed when specifying which
$U(1)$ subgroup is factored.
\end{remark}

%──────────────────────────────────────────────────────────────────────────────
\section{Exterior algebra and the Fock representation}
\label{sec:exterior}
%──────────────────────────────────────────────────────────────────────────────

\subsection{Construction of $\Lambda^\bullet V$}

Since $\dim_\mathbb{C} V = 3$, the exterior algebra is
\begin{equation}
  \Lambda^\bullet V = \Lambda^0 V \oplus \Lambda^1 V \oplus \Lambda^2 V \oplus \Lambda^3 V,
  \qquad \dim_\mathbb{C} = 1+3+3+1 = 8.
\end{equation}
Basis elements, with occupation-number labels (three fermionic modes):
\begin{alignat}{3}
  \Lambda^0 V &:\quad \lvert 000\rangle \leftrightarrow 1, & \quad & \text{(vacuum)}\\
  \Lambda^1 V &:\quad \lvert 100\rangle \leftrightarrow I,\quad
                       \lvert 010\rangle \leftrightarrow J,\quad
                       \lvert 001\rangle \leftrightarrow K, &&\\
  \Lambda^2 V &:\quad \lvert 110\rangle \leftrightarrow I\wedge J,\quad
                       \lvert 101\rangle \leftrightarrow I\wedge K,\quad
                       \lvert 011\rangle \leftrightarrow J\wedge K, &&\\
  \Lambda^3 V &:\quad \lvert 111\rangle \leftrightarrow I\wedge J\wedge K.
                      &&\text{(full)}
\end{alignat}

\emph{Note.}  The binary occupation-number labels are a computational
representation for three fermionic modes.  They do not assert that physical
spacetime is literally a three-qubit quantum computer.  The theorem holds for
any choice of Jordan--Wigner ordering convention.

\subsection{Decomposition under SU(3)}

The $\mathrm{SU}(3)$ action on $V$ induces an action on each $\Lambda^k V$.

\begin{theorem}\label{thm:fock_reps}
Under the induced $\mathrm{SU}(3)$ action:
\begin{enumerate}
  \item $\Lambda^0 V \cong \mathbf{1}$ (trivial representation);
  \item $\Lambda^1 V \cong \mathbf{3}$ (fundamental representation);
  \item $\Lambda^2 V \cong \bar{\mathbf{3}}$ (anti-fundamental representation);
  \item $\Lambda^3 V \cong \mathbf{1}$ (trivial representation, signed by $\det g = 1$).
\end{enumerate}
Hence $\Lambda^\bullet V \cong \mathbf{1}\oplus\mathbf{3}\oplus\bar{\mathbf{3}}\oplus\mathbf{1}$.
\end{theorem}

\begin{proof}
\emph{(1)} $\Lambda^0 V = \mathbb{C}$ with trivial action.

\emph{(2)} $\Lambda^1 V = V$ with the defining (fundamental) action of $\mathrm{SU}(3)$.

\emph{(3)} $\Lambda^2 V$.
For $\mathrm{SU}(3)$, the two-index antisymmetric tensor in the fundamental
representation is the anti-fundamental $\bar{\mathbf{3}}$.  Explicitly,
the isomorphism $\Lambda^2 V \cong V^*$ is given by contraction with the invariant
$\epsilon_{ijk}$ tensor:
\begin{equation}
  (e_i\wedge e_j) \mapsto \epsilon_{ijk}\,e^k \in V^*,
\end{equation}
where $(e_1,e_2,e_3) = (I,J,K)$ and $\epsilon_{ijk}$ is the Levi-Civita symbol.
Since $\mathrm{SU}(3)$ preserves $\epsilon_{ijk}$ (as $\det g = 1$), this map is
$\mathrm{SU}(3)$-equivariant, giving $\Lambda^2 V \cong V^* \cong \bar{\mathbf{3}}$.

\emph{(4)} $\Lambda^3 V = \mathbb{C}\cdot(I\wedge J\wedge K)$ is one-dimensional.
For $g\in\mathrm{SU}(3)$, $g(I\wedge J\wedge K) = (\det g)(I\wedge J\wedge K) = I\wedge J\wedge K$.
Hence the action is trivial.
\end{proof}

\begin{remark}[Relation to earlier ``vector/bivector'' intuition]
The identification
\begin{equation}
  \text{vector sector } \Lambda^1 V = \mathbf{3}, \qquad
  \text{bivector sector } \Lambda^2 V = \bar{\mathbf{3}}
\end{equation}
gives a precise version of the earlier intuition about vector and bivector
subspaces.  It is distinct from the eight-dimensional adjoint representation
$\mathbf{8}$, which consists of the generators of $\mathfrak{su}(3)$.
\end{remark}

%──────────────────────────────────────────────────────────────────────────────
\section{Second-quantised generators on $\Lambda^\bullet V$}
\label{sec:fock}
%──────────────────────────────────────────────────────────────────────────────

\subsection{The 8 Fock states are not 8 gluons}

\begin{remark}\label{rem:states_not_gluons}
The $8$-dimensional exterior space may be represented as a three-mode
fermionic Fock space with occupation basis
$\{|000\rangle,|100\rangle,\ldots,|111\rangle\}$.  This is a
representation-theoretic construction; UBT has not thereby derived three
fundamental physical fermionic modes.  The eight QCD gauge-boson species are
labelled by the eight Lie-algebra directions $T_a$ (the adjoint index $a$), not
by these eight Fock basis states.  The operators $\hat T_a$ below are the
induced representation of those generators on $\Lambda^\bullet V$; the
$8$-dimensional Fock representation is not the adjoint $\mathbf 8$.
\end{remark}

\subsection{Fermionic creation/annihilation operators}

Let $(e_1,e_2,e_3) = (I,J,K)$.  Define fermionic creation operators
$c_i^\dagger$ and annihilation operators $c_i$ on $\Lambda^\bullet V$ by
\begin{align}
  c_i^\dagger : \Lambda^k V &\to \Lambda^{k+1} V, \quad
    c_i^\dagger (v) = e_i \wedge v,\\
  c_i : \Lambda^k V &\to \Lambda^{k-1} V, \quad
    c_i = (c_i^\dagger)^*.
\end{align}
These satisfy the canonical anticommutation relations (CAR):
\begin{equation}
  \{c_i, c_j^\dagger\} = \delta_{ij},\quad
  \{c_i, c_j\} = 0,\quad
  \{c_i^\dagger, c_j^\dagger\} = 0.
\end{equation}

\subsection{Second-quantised su(3) generators}

Let $T_a = \lambda_a/2$ for $a=1,\ldots,8$ be the $\mathfrak{su}(3)$ generators
in the fundamental representation (with the standard Gell-Mann matrices
$\lambda_a$ satisfying $\mathrm{tr}(\lambda_a\lambda_b) = 2\delta_{ab}$).
Define the second-quantised operators
\begin{equation}\label{eq:Thatdef}
  \hat{T}_a := \sum_{i,j=1}^3 c_i^\dagger (T_a)_{ij} c_j.
\end{equation}

\begin{theorem}\label{thm:su3_algebra}
The operators $\hat{T}_a$ satisfy
\begin{equation}
  [\hat{T}_a, \hat{T}_b] = i f_{abc}\,\hat{T}_c,
\end{equation}
where $f_{abc}$ are the $\mathfrak{su}(3)$ structure constants.
\end{theorem}

\begin{proof}
For number-preserving fermionic bilinears, the general identity
$[\mathrm{d}\Gamma(A), \mathrm{d}\Gamma(B)] = \mathrm{d}\Gamma([A,B])$
holds, where $\mathrm{d}\Gamma(A) = \sum_{ij} c_i^\dagger A_{ij} c_j$
is the second-quantised lift of $A\in\mathrm{End}(V)$.  Since
$[T_a, T_b] = if_{abc} T_c$, we obtain
$[\hat{T}_a, \hat{T}_b] = \mathrm{d}\Gamma([T_a,T_b]) = if_{abc}\hat{T}_c$.
\end{proof}

\subsection{Block structure on $\Lambda^\bullet V$}

Under particle number, the action of $\hat{T}_a$ is block-diagonal:
\begin{equation}
  \hat{T}_a = 0 \oplus (T_a)^{[1]} \oplus (T_a)^{[2]} \oplus 0,
\end{equation}
where the blocks correspond to $1+3+3+1$ and
$(T_a)^{[1]}$ is the fundamental action on $\Lambda^1 V$,
$(T_a)^{[2]}$ is the anti-fundamental action on $\Lambda^2 V$.

The number operator $\hat{N} = \sum_i c_i^\dagger c_i$ commutes with all
$\hat{T}_a$ since $\mathrm{d}\Gamma$ preserves particle number.

%──────────────────────────────────────────────────────────────────────────────
\section{Machine verification}
\label{sec:verification}
%──────────────────────────────────────────────────────────────────────────────

The theorem-critical algebraic identities listed below are verified by the exact symbolic script
\begin{center}
  \texttt{verification/su3\_stabilizer\_fock\_check.py}
\end{center}
using SymPy (exact integer/rational arithmetic, no floating-point tolerances).
The script checks:
\begin{enumerate}
  \item Quaternion multiplication identities and the canonical Hermitian form
        $h(v,w)=\mathrm{Sc}(v^\dagger w)$;
  \item The symbolic determinant identity for $\Omega$ (Theorem~\ref{thm:Omega});
  \item Gell-Mann matrices are Hermitian and traceless;
  \item All 28 independent commutator pairs reproduce the standard $\mathfrak{su}(3)$ structure constants;
  \item Explicit $8\times 8$ fermionic creation/annihilation matrices via
        Jordan--Wigner convention satisfy CAR exactly;
  \item All $8$ second-quantised operators $\hat{T}_a$ are constructed;
  \item All $28$ commutators of $\hat{T}_a$ are verified;
  \item The number operator commutes with all $\hat{T}_a$;
  \item The particle-number blocks have dimensions $1,3,3,1$;
  \item The one-particle block is the fundamental $\mathbf{3}$;
  \item The two-particle block is equivalent to the anti-fundamental $\bar{\mathbf{3}}$;
  \item Vacuum and fully occupied sectors are singlets.
\end{enumerate}
Run with: \texttt{python verification/su3\_stabilizer\_fock\_check.py}

\medskip
\textbf{Lean formalisation:} \texttt{LEAN-PENDING}.
Priority items: stabiliser implication (unitary + det=1 $\Rightarrow$ SU(3)),
exterior-dimension identities, and the abstract $\mathrm{d}\Gamma$ commutator lemma.

%──────────────────────────────────────────────────────────────────────────────
\section{Dynamical gauge gap}
\label{sec:dyn_gap}
%──────────────────────────────────────────────────────────────────────────────

\begin{gapbox}
\textbf{GAP-SU3-DYN: OPEN}

The results above establish the algebraic $\mathrm{SU}(3)$ stabiliser of the
canonical biquaternionic colour sector.  The following additional steps are
required to identify this with the dynamical local QCD gauge symmetry, and
none of them has yet been derived from the canonical UBT action:

\begin{itemize}
  \item Does the canonical UBT action locally preserve \emph{both} $h$ and
        $\Omega$ on the colour sector?
  \item Can the local frame redundancy for $(h,\Omega)$ be derived from the
        UBT covariant derivative to give a local $\mathrm{SU}(3)$ gauge connection?
  \item Does the Yang--Mills kinetic term for the colour connection arise from
        the canonical UBT action?
\end{itemize}

Until these questions are answered, the separation is:

\begin{description}
  \item[PROVED:] The canonical biquaternionic colour sector carries a natural
    $\mathrm{SU}(3)$ stabiliser and the exterior/Fock representation
    $\mathbf{1}\oplus\mathbf{3}\oplus\bar{\mathbf{3}}\oplus\mathbf{1}$.
  \item[OPEN/CONDITIONAL:] Identifying this stabiliser with the dynamical
    local QCD gauge symmetry requires an action-level/local-frame derivation.
  \item[OPEN:] Dynamical confinement / mass gap (Clay Millennium problem).
\end{description}
\end{gapbox}

\ifdefined\INCLUDEMODE
\else
\end{document}
\fi
